\documentclass[a4paper]{article}
\usepackage{amsfonts}
\usepackage{amsmath}
\usepackage{amssymb}
\usepackage{graphicx}     

\begin{document}
  
\noindent \large Auteur: Michael Zielinski \\
\noindent \large Studierichting: Informatica\\
\noindent \large Oefeningenreeks 6F nr. 3\\

\medskip

\normalsize

$ f_n : \mathbb{R}_0^+ \rightarrow \mathbb{R} : x \mapsto \dfrac{1}{nx^2} $ \\

We bepalen eerst de puntsgewijze limiet: \\

$ \lim\limits_{n \rightarrow \infty} f_n(x)
= \lim\limits_{n \rightarrow \infty}\dfrac{1}{nx^2} $ \\

Voor elke vaste $x \in \mathbb{R}_0^+$ geldt: \\

$ \lim\limits_{n \rightarrow \infty}\dfrac{1}{nx^2} = 0 $ \\

Dus: \\

$ f : \mathbb{R}_0^+ \rightarrow \mathbb{R} : x \mapsto 0 $ \\

$ f_n \stackrel{p}{\rightarrow} f $ \\

Nu onderzoeken we uniforme convergentie: \\

$ \left|\left|f_n-f\right|\right|_{\infty}
= \sup\limits_{x \in \mathbb{R}_0^+}\left|\dfrac{1}{nx^2}-0\right| $ \\

$ = \sup\limits_{x \in \mathbb{R}_0^+}\dfrac{1}{nx^2} $ \\

Omdat $x$ willekeurig dicht bij $0$ kan liggen, wordt $\dfrac{1}{nx^2}$ onbeperkt groot. \\

Dus: \\

$ \left|\left|f_n-f\right|\right|_{\infty}=+\infty $ \\

Hieruit volgt: \\

$ \lim\limits_{n \rightarrow \infty}\left|\left|f_n-f\right|\right|_{\infty}\neq 0 $ \\

Dus: \\

$ f_n \stackrel{u}{\not\rightarrow} f $ \\

\begin{thebibliography}{999}
%\bibitem{a} Referentie
\end{thebibliography}

\end{document}